\documentclass[10pt,a4paper]{amsart}
\usepackage[margin=19mm]{geometry}
\usepackage{amsmath,amssymb,mathtools}
\usepackage[scr=rsfs,cal=euler]{mathalfa}
\usepackage{xfrac}
\usepackage[unicode,hidelinks,bookmarks=false]{hyperref}
\newcommand{\ie}{i.e.\ }
\newcommand{\cf}{cf.\ }
\newcommand{\resp}{respectively\ }
\newcommand{\Subsection}{\subsection}
\providecommand{\nobreakdash}{\nobreak\hbox{-}}
\setlength{\emergencystretch}{2em}
\multlinegap0pt
\usepackage{xfrac}
\usepackage{amssymb}
\usepackage{mathtools}
\newtheorem{theorem}{Theorem}
\newtheorem{lemma}{Lemma}
\DeclareMathOperator{\Max}{Max}
\newcommand{\frakL}{\mathfrak L}
\newcommand{\frakM}{\mathfrak M}
\newcommand{\frakS}{\mathfrak S}
\title{Remarks on infimum and maximal lower bounds\\ of a set of bounded self-adjoint operators}
\author{Matthias G\"unther, Lutz Klotz}
\date{Source: arXiv:2604.22488v1 (CC BY 4.0)}
\begin{document}
\maketitle

\noindent\emph{Source and licence.} Remarks on infimum and maximal lower bounds\\ of a set of bounded self-adjoint operators. Version arXiv:2604.22488v1 (CC BY 4.0), distributed by arXiv under the Creative Commons Attribution 4.0 International licence, http://creativecommons.org/licenses/by/4.0/, as recorded in the arXiv licence metadata for this exact version and shown on the abstract page https://arxiv.org/abs/2604.22488v1. Full licence notice: \texttt{licenses/arxiv-2604.22488-CC-BY-4.0.txt}.\par\smallskip

\noindent\textit{Editorial scope.} Counted unit: the article's theorem thm5.5 (source0.tex lines 1636--1644). The article's complete original proof of it is delivered with it (source0.tex lines 1646--1785, one \texttt{proof} environment of 140 lines). No mathematics is written, repaired or re-derived: the statement and the proof are the article's own lines, in the article's own notation, and no retained line is substituted or re-flowed.  Retained uncounted as the source-local context the proof needs: lines 690--703; lines 1008--1011; lines 1576--1588.  Nothing was omitted from any retained span, so the package carries no omission record.  No retained line contains a citation, so the wrapper carries no bibliography and no bibliography entry.  No retained line contains a figure environment, an image-inclusion command or drawing code, so no image is delivered and the package declares no asset dependency.  Editorial formatting only, in addition: an AMS \texttt{amsart} preamble that restates the article's own declarations and macros used by the retained lines, namely the environments \texttt{theorem}, \texttt{lemma} on separate counters, together with the standard packages those lines need.  The retained environments print in this document as Theorem 1, Lemma 1 in order of appearance, whereas the article prints the same items with the numbering of its own full document.  The complete native source of the article as distributed in the arXiv e-print for this exact version is bundled unchanged as \texttt{sources/199098/source0.tex}.
\begin{theorem}
  \label{thm2.2}
An operator $S\in\frakS$ is positive if and only if the following three
conditions are satisfied:
\begin{enumerate}
\item[\normalfont(i)] $S_{1}\geq0$,  
\item[(ii)] $\mathscr R(S_{12})\subseteq\mathscr R(S_{1}^{\sfrac{1}{2}})$,
\item[(iii)]
$S_{2}-\bigl((S_{1}^{\#})^{\sfrac{1}{2}}S_{12}\bigr)^{\ast}
(S_{1}^{\#})^{\sfrac{1}{2}}S_{12}\geq0$.
\end{enumerate}
Here, $(S_{1}^{\#})^{\sfrac{1}{2}}$ denotes the unique possibly unbounded
positive self-adjoint square root of $S_{1}^{\#}$.
\end{theorem}

\begin{equation}
\label{eq4.3}
\overline{\sum_{A\in\frakM}\mathscr N(A-M)} = \mathscr H,  
\end{equation}

\begin{lemma}
\label{lem5.4}
Let $\mathscr K$ be a subspace of $\mathscr H$.
\begin{enumerate}
\item[(i)]
  If there exists $M\in\Max\frakL$ satisfying $Mu=Au$, $u\in\mathscr
  K$, then $(Au,u)<0$ for all $u\in\mathscr K\setminus\{0\}$.
\item[(ii)]
  If there exists $M\in\Max\frakL$ such that $\mathscr
  K\subseteq\mathscr N(M)$, then $(Au,u)>0$ for all  $u\in\mathscr
  K\setminus\{0\}$.
\end{enumerate}
\end{lemma}

\begin{theorem}
\label{thm5.5}
The map
\[
X\mapsto M(X) := J - S(X), \; X\in\mathbb C^{p\times q}
\]
establishes a one-to-one correspondence between
$\mathbb C^{p\times q}$ and $\Max\frakL$. 
\end{theorem}

\begin{proof}
Step 1: For $X\in\mathbb C^{p\times q}$, the matrix $M(X)$ is a
maximal lower bound of $\frakM$. Since $S(X)\in\frakS^{+}$ by
Theorem~\ref{thm2.2}, $M(X)\leq J$. Since
\[
M(X)=
- \begin{pmatrix}
  XX^{\ast} & (I_{p}+XX^{\ast})^{\sfrac{1}{2}}X \\
  X^{\ast}(I_{p}+XX^{\ast})^{\sfrac{1}{2}} & I_{q}+X^{\ast}X
  \end{pmatrix}
\]
and
\[
(I_{p}+XX^{\ast})^{\sfrac{1}{2}}X = X(I_{q}+X^{\ast}X)^{\sfrac{1}{2}}
\]
another application of Theorem~\ref{thm2.2} yields $M(X)\leq0$. To
prove that $M(X)$ is a maximal element of $\frakL$ note first that
\[
  \mathscr N\bigl(S(X)\bigr)=\left\{
    \begin{pmatrix}
    -(I_{p}+XX^{\ast})^{-\sfrac{1}{2}} X x_{2} \\ x_{2}
    \end{pmatrix}
    : x_{2} \in\mathscr H_{2}\right\},
\]  
which yields
\[
\dim\mathscr N\bigl(J-M(X)\bigr) = \dim\mathscr N\bigl(S(X)\bigr)\geq q.
\]
Similarly,
\[
  \mathscr N\bigl(M(X)\bigr)=\left\{
    \begin{pmatrix}
    x_{1} \\
    -(I_{q}+X^{\ast}X)^{-\sfrac{1}{2}} X^{\ast}x_{1} 
    \end{pmatrix}
    : x_{1} \in\mathscr H_{1}\right\},
\]
hence, $\dim\mathscr N\bigl(M(X)\bigr)\geq p$.
Since $M(X)+S(X)=J$ yields
\[
\mathscr N\bigl(M(X)\bigr) \cap \mathscr N\bigl(S(X)\bigr) = \{0\},
\]
we obtain 
\[
\mathscr N\bigl(0-M(X)\bigr) + \mathscr N\bigl(J-M(X)\bigr) = \mathbb C^{n},
\]
and $M(X)\in\Max\frakL$ by \eqref{eq4.3}. \\
Step~2: The correspondence $X\mapsto M(X)$ is one-to-one. Let
$X,Y\in\mathbb C^{p\times q}$ and $M(X)=M(Y)$. Comparing the left
upper  corners of $M(X)$ and $M(Y)$ we get
$XX^{\ast}=YY^{\ast}$. Since $I_{p}+XX^{\ast}$ is invertible, it
follows $X=Y$ from the equality of the right upper corners of $M(X)$
and $M(Y)$. \\
Step~3: For arbitrary $M\in\Max\frakL$, there exists $X\in\mathbb
C^{p\times q}$ satisfying $M=M(X)$. By Lemma~\ref{lem5.4} the space
$\mathscr N(J-M)$ is a negative definite subspace and the space
$\mathscr N(0-M)$ is a positive definite subspace with respect to the
indefinite metric generated by $J$. Since 
\begin{equation}
\label{eq5.4}  
\mathscr N\bigl(0-M\bigr) + \mathscr N\bigl(J-M\bigr) = \mathscr H
\end{equation}
by \eqref{eq4.3}, the space $\mathscr N(0-M)$ is a maximal positive
definite subspace. Let $K$ be its angular operator. Recall that
\[
  \mathscr N(0-M) = \left\{
     \begin{pmatrix} x_{1} \\ Kx_{1} \end{pmatrix}
      : x_{1}\in\mathscr H_{1} \right\}
\]  
and $K$ is a contraction from $\mathscr H_{1}$ into $\mathscr H_{2}$
with the additional property that $I_{p}-K^{\ast}K$ is invertible.
Setting
\[
X:=-(I_{p}-K^{\ast}K)^{-\sfrac{1}{2}}K^{\ast},
\]
we obtain
\begin{align*}
(I_{p}+XX^{\ast}) & =
  I_{p}+(I_{p}-K^{\ast}K)^{-\sfrac{1}{2}}K^{\ast}K
  (I_{p}-K^{\ast}K)^{-\sfrac{1}{2}} 
   = I_{p} + (I_{p}-K^{\ast}K)^{-1} K^{\ast}K \\                    
 & = (I_{p}-K^{\ast}K)^{-1}(I_{p}-K^{\ast}K + K^{\ast}K)
   = (I_{p}-K^{\ast}K)^{-1}
\end{align*}  
and 
\[
X^{\ast}X = K(I_{p}-K^{\ast}K)^{-1}K^{\ast},
\]
hence   
\[
  S(X)=
  \begin{pmatrix}
    (I_{p}-K^{\ast}K)^{-1}    & -(I_{p}-K^{\ast}K)^{-1}K^{\ast} \\
    -K (I_{p}-K^{\ast}K)^{-1} & K (I_{p}-K^{\ast}K)^{-1} K^{\ast}
  \end{pmatrix}.
\]
Now it is not hard to see that for 
\[
u = \begin{pmatrix} x_{1} \\ K x_{1} \end{pmatrix}
    \in\mathscr N(0-M),\quad
M(X)u = \bigl(J-S(X)\bigr) \begin{pmatrix} x_{1} \\ K x_{1} \end{pmatrix}=0,
\]  
thus
\begin{equation}
\label{eq5.5} 
  Mu = M(X)u, u\in\mathscr N(0-M).
\end{equation}
If $u\in\mathscr N(J-M)$ and $v\in\mathscr N(M)$, we have
$(Ju,v)=(Mu,v)=(u,Mv)=0$. It follows that $\mathscr N(J-M)$ is a
subspace of the $J$-orthogonal complement of $\mathscr N(M)$. Since
the $J$-orthogonal complement of a maximal positive definite subspace
$\mathscr N(M)$ is a negative definite subspace and since $\mathscr
N(J-M)$ is a maximal negative definite subspace, we can conclude
that $\mathscr N(J-M)$ is the $J$-orthogonal complement of
\[
 \mathscr N(M) =
 \left\{ \begin{pmatrix} x_{1}\\ K x_{1} \end{pmatrix} :
  x_{1}\in\mathscr H_{1} \right\},
\]  
which yields
\[
  \mathscr N(J-M)=
  \left\{ \begin{pmatrix} K^{\ast}x_{2}\\x_{2} \end{pmatrix} :
  x_{2}\in\mathscr H_{2} \right\}.
\]  
It is easy to see that
\[
  M(X)\begin{pmatrix} K^{\ast}x_{2}\\x_{2} \end{pmatrix}
  = \bigl(J-S(X)\bigr)
    \begin{pmatrix} K^{\ast}x_{2}\\x_{2} \end{pmatrix}
  = \begin{pmatrix} K^{\ast}x_{2}\\-x_{2} \end{pmatrix}
\] 
or, equivalently, $M(X)v=Jv$, $v\in\mathscr N(J-M)$. It follows
\begin{equation}
\label{eq5.6}
Mv = M(X)v,\;v\in\mathscr(J-M).  
\end{equation}
Finally, equality $M=M(X)$ is a consequence of \eqref{eq5.4},
\eqref{eq5.5}, and \eqref{eq5.6}.
\end{proof}  

\end{document}
